\begin{pastexam}[label={px:nccu-110-dsa-04}]{政大資科所 110 學年度 第 4 題}
（問答）

(10\%) Give two advantages of red black trees over hash tables.
\end{pastexam}

\begin{pastsoln}
下面 $n$ 是存的鍵數。 紅黑樹的高度不超過 $2\lg(n+1)$（資結書定理~10.4）。

\textbf{優點一：保有鍵的順序，支援和順序有關的操作}。 紅黑樹是二元搜尋樹，中序走訪就是由小到大。 找最小、最大值，
找某個鍵的前一個或後一個（successor），都是 $O(\log n)$；列出某個範圍 $[x, y]$ 裡的鍵是 $O(\log n + k)$（$k$ 是
範圍裡的鍵數）；由小到大輸出全部是 $O(n)$。 雜湊表把鍵打散存放，沒有順序，這些操作都要掃過整張表，$\Theta(n + m)$
（$m$ 是表的大小），排序輸出還要再 $O(n \log n)$。 例如「找出學號在 $1100$ 到 $1199$ 之間的學生」，紅黑樹很方便，
雜湊表做不到快。

\textbf{優點二：最壞情況的保證}。 紅黑樹的搜尋、插入、刪除最壞都是 $O(\log n)$，不依賴輸入的分布。 雜湊表的 $O(1)$
是\textbf{期望值}，前提是雜湊函數把鍵打得夠散（簡單均勻雜湊）；最壞情況所有鍵撞在同一格，每次操作 $\Theta(n)$。
惡意挑選的輸入可以故意製造碰撞。 另外雜湊表放太滿要整張重新雜湊，那一次插入要 $\Theta(n)$（平攤才是常數），
紅黑樹每次插入、刪除最多做常數次旋轉，沒有這種突然很慢的操作。

其他也可以寫的：紅黑樹的空間和 $n$ 成正比，不必預先決定表的大小或挑雜湊函數；鍵只要能比大小就行，不需要設計
雜湊函數。 （反過來，雜湊表的優點見歷屆試題~\ref{px:nccu-108-dsa-02}；資結書第~10.5~節有兩者的比較。）
\end{pastsoln}
